\exercisegroup

\begin{enumerate}
\def\labelenumi{\arabic{enumi}.}
\tightlist
\item
  证明关系
\end{enumerate}

\[
(x = y) \iff (\forall X) ((x \in X) \Longrightarrow (y \in X))
\]

是一个定理。

\begin{enumerate}
\def\labelenumi{\arabic{enumi}.}
\setcounter{enumi}{1}
\item
  证明 \(\phi \neq \{x\}\) 是一个定理。由此推出 \((\exists x)(\exists y)(x \neq y)\) 是一个定理。
\item
  设 A 和 B 是集合 X 的两个子集。证明关系 \(B \subset GB\) 等价于 \(A \subset GB\) 并且关系 \(GB \subset A\) 等价于 \(CA \subset B\) .
\item
  证明关系 \(\mathbf{X} \subset \{x\}\) 等价于
\end{enumerate}

\[
" \mathrm{X} = \{x \} \text { 或 } \mathrm{X} = \phi^ {\prime \prime}.
\]

\begin{enumerate}
\def\labelenumi{\arabic{enumi}.}
\setcounter{enumi}{4}
\item
  证明 \(\phi = \tau_{\mathbf{X}}(\tau_x(x \in \mathbf{X}) \notin \mathbf{X})\) .
\item
  让 \(\mathfrak{C}\) 是一个包含该符号的平等主义理论 \(\in\) 以及以下公理：
\end{enumerate}

A1'。

\[
(\forall y) (y = \tau_ {x} ((\forall z) (z \in x \Longleftrightarrow z \in y)))
\]

（换言之，每个项都等于其元素的集合）。证明外延公理A1在 \(\mathfrak{G}\) 中是定理（使用模式S7）。

